\chapter{The marginal problem}
\label{ch:marginal}

% Part V, added 2026-09-25 from ~/Downloads/chygraph_master_equation/draft.tex,
% Sec. 9 (category theory and the marginal problem). Nothing is computed here.
% The chapter attributes Ch. 17's exactness condition to the marginal problem
% of the 1960s, reads Part IV's error as a cohomological obstruction, and
% states the level-wise gluing theorem as a corollary of Vorob'ev-Kellerer
% modulo one definition check (the two notions of consistency coincide).
% Revised 2026-09-25 after the first Part V review (items A1-A3, A5, B; the
% review file is gone, extensions.md records what changed) and 2026-09-26
% after ~/Downloads/book_review.md (A7, B23).

Chapter~\ref{ch:chyequation} put every calculation in the book on one
bipartite graph, the factor graph of $\alpha+I$, and said that the recursion
is exact when that graph is a tree, which is Berge-acyclicity of the
complex family. Chapter~\ref{ch:exactness} had found a weaker condition
from the other side, by measurement, for the region-graph calculation of
Chapter~\ref{ch:gbp}, and given it three names, of which
$\alpha$-acyclicity is the one this chapter needs. This chapter says where
that condition comes from. It is not a fact
about belief propagation. It is a fact about probability measures on a
hypergraph --- when a family of marginals, one per hyperedge, is the shadow of
a single joint law --- and it was settled in the early 1960s, twenty years
before the algorithm the book runs and sixty before this book.

The chapter has two halves. The first is the categorical reading of the
recursion itself: the sum-product update as a computation in a monoidal
category, the treelike assumption as a statement about conditional
independence, which in the categories built for the purpose is a definition
rather than a heuristic. The second is the marginal problem proper, its
answer, the sheaf-theoretic account of what happens when the answer is no,
and the one question the chygraph adds: marginals of marginals.

\section{The recursion in a category}
\label{sec:mg-category}
\index{monoidal category}
\index{hypergraph category}
\index{Markov category}
\index{conditional independence|textbf}

\paragraph{The connection.} \citet{morton2014} runs the sum-product
algorithm in any monoidal category equipped with a copy-and-delete
structure: objects are the variables, morphisms are the factors, copying a
wire is what lets one variable enter several factors and deleting one is what
summing it out means. The algorithm is stated once for the category and
specialises to belief propagation on a factor graph when the category is that
of finite sets with non-negative kernels between them. Chygraph belief
propagation, Eqs.~\eqref{eq:chyup}--\eqref{eq:chydown}, is an instance of that
statement with the factor graph of rule~3 of Sec.~\ref{sec:ce-representation}
as the diagram: every vertex a wire, every complex a box with one leg per
member and one leg for its own state, Fig.~\ref{fig:factorgraph} read as a
string diagram. Nothing in \citeauthor{morton2014}'s construction has to
change for a box whose output wire enters another box, which is nesting, and
that is the categorical way of saying what Sec.~\ref{sec:ce-known} said: the
same object is iterated at every level.

The \emph{hypergraph categories} of \citet{fong2019} make the wiring itself
algebraic. A variable is a Frobenius spider --- a node that any number of
wires may enter and leave, with the copy and delete maps as its structure
--- and a factor is a morphism between the spiders it touches. A diagram in
such a category is exactly a factor graph, and a diagram that is a tree
composes sequentially: it can be evaluated by contracting one box at a time
from the leaves inward, with each contraction a morphism composition and no
step ever needing to revisit a wire. That is the exactness of belief
propagation on a tree, drawn rather than proved, and the drawing is
Fig.~\ref{fig:twosteps} of Chapter~\ref{ch:statmech} with its two steps read
as the two kinds of composition a tree admits: composing along a wire, which
is the up step, and composing into a box, which is the down step.

\paragraph{Treelike as conditional independence.}
\citet{fritz2020}'s Markov categories are the synthetic setting for
conditional independence: a category in which ``$X$ and $Y$ are independent
given $Z$'' is a statement about how a morphism factors through copies of
$Z$, provable from axioms rather than checked on densities. In that language
the assumption the book makes at Eq.~\eqref{eq:up}, and names in
Secs.~\ref{sec:treelike} and~\ref{sec:assumes}, is one sentence: \emph{the
complexes at an atom carry independent information given the atom}. Adding
the fields $u_{b\to i}$ independently is exact if and only if the
messages arriving at $i$ from its different containers are conditionally
independent given $\sigma_{i}$, and that is a factorisation of a morphism,
not a property of a solution. Section~\ref{sec:price} said what it costs when
false --- two complexes sharing two atoms make the two messages at either
shared atom depend on one another through the other --- and the Markov
category says the same thing without a model: the diagram has a cycle, so the
morphism does not factor.

The value of the categorical statement is not that it proves anything the
book does not already know. It is that it separates the assumption from the
model. Every row of Table~\ref{tab:ce-substitutions} makes the same
conditional-independence assumption, whatever it iterates, and the assumption
is a property of the diagram alone.

\section{The marginal problem}
\label{sec:mg-marginal}
\index{marginal problem|textbf}
\index{acyclic hypergraph}
\index{join tree}
\index{Vorob'ev's theorem|textbf}
\index{chordal graph}

\paragraph{The question.} Take a hypergraph $H$ on a set of variables and,
for every hyperedge $e$, a probability law $\mu_{e}$ on the variables in
$e$. Say the family is \emph{consistent} if any two of its members agree on
their overlap: for every pair $e,e'$ the marginal of $\mu_{e}$ on $e\cap e'$
equals that of $\mu_{e'}$. Does a consistent family always come from one
joint law $\mu$ on all the variables, of which each $\mu_{e}$ is the
marginal?

For a tree of pairwise hyperedges, yes: build $\mu$ outward from a root,
attaching each new variable by its conditional law given its parent, and
consistency guarantees that every edge marginal is reproduced. For a
triangle of pairwise hyperedges, not always: take three binary variables with
uniform marginals, let the laws on $\{1,2\}$ and $\{2,3\}$ be perfectly
correlated and the law on $\{1,3\}$ perfectly anticorrelated; the three
agree on every overlap, and no joint law has all three as marginals. The general answer is due to
\citet{vorobev1962} and \citet{kellerer1964}, working independently:

\begin{center}
\emph{every consistent family of marginals on $H$ glues to a global law\\
if and only if $H$ is acyclic.}
\end{center}

The ``if'' is the tree construction with a hyperedge in place of an edge:
order the hyperedges so that each meets the union of its predecessors in a
subset of one of them, and attach. The ``only if'' is the triangle
generalised: on any cycle a consistent family can be built with no global
extension. What \citeauthor{vorobev1962} called acyclic was, in the
vocabulary that came later, the condition that the hyperedges admit such an
ordering, and \citet{beeri1983} and \citet{fagin1983} identified it: it is
$\alpha$-acyclicity, the weakest of their notions, and it is equivalent to
the existence of a join tree and to the GYO reduction of \citet{graham1979}
and \citet{yu1979} consuming the whole family. The ordering in the ``if''
direction is the reverse of a GYO run.

\paragraph{In the book's vocabulary.} Section~\ref{sec:threenames} stated
the chain chordal $=$ join tree $=$ $\alpha$-acyclic for the maximal-clique
family of a graph and certified it by leaf removal on the incidence
structure. That chain is the Vorob'ev--Kellerer theorem's hypothesis, and
the theorem is its conclusion in the form the book needs, with one care
about which fixed point it applies to. The theorem's consistency is
agreement on the \emph{whole} overlap $e\cap e'$ of two hyperedges. The
messages of Eqs.~\eqref{eq:chyup}--\eqref{eq:chydown} enforce agreement on
single atoms only --- the marginal at atom $i$ is the same whichever
complex it is read from --- so on two triangles sharing the pair
$\{2,3\}$ the two beliefs need not agree on the pair, and the chygraph
fixed point is not a consistent family in the theorem's sense. The
region-graph fixed point of Chapter~\ref{ch:gbp} is: closing the family
under intersection, rule~2 of Sec.~\ref{sec:regions}, makes every overlap a
region with a belief of its own, and the parent-to-child messages of
Eq.~\eqref{eq:p2c} enforce agreement on it. Both statements are measured
(Sec.~\ref{sec:mg-checks}). On the $240$ runs of Chapter~\ref{ch:overlap}
at the stable fixed point, the beliefs of two cliques sharing a pair of
atoms agree on each atom to $10^{-13}$ and differ on the pair by up to
$0.15$ with every bond in every clique and $0.33$ with each bond in one,
on $101$ and $200$ of the $240$ by more than $10^{-6}$; the disagreement
lives at the weaker coupling, since at the stronger one every belief is
concentrated on the ordered configuration. On the $186$ runs of
Chapter~\ref{ch:gbp} whose region-graph iteration reached a fixed point,
the beliefs agree on every overlap to $10^{-8}$ or better, chordal or not:
on the $51$ non-chordal ones among them the family is consistent and is
not the family of true marginals, its $\ln Z$ being off by up to $12$. That family, when the
hypergraph is $\alpha$-acyclic, glues to a joint law. That the joint law is
the Boltzmann measure, and Eq.~\eqref{eq:bethe}'s free energy exact, is a
second step: it needs the junction-tree factorisation of the joint law as
a product of clique marginals over a product of separator marginals,
together with \citeauthor{lauritzen1988}'s result that the fixed point on a
join tree is the true marginals. When the hypergraph is not
$\alpha$-acyclic there exists a consistent family with no extension, and
nothing prevents the recursion from converging to one.

So the exactness theorem of Chapter~\ref{ch:exactness} is sixty years older
than belief propagation, and older than the region-graph repair by the same
margin. \citet{lauritzen1988} proved that propagation over a join tree is
exact, which is the algorithmic form; the probabilistic form, that a join
tree is precisely when there is something exact to compute, is
\citeauthor{vorobev1962}'s. Chapter~\ref{ch:gbp} measured it without naming
it: Sec.~\ref{sec:gbpverdict} found the region-graph fixed point exact on
every chordal instance and on no other over $240$ runs, which is consistent
with both directions of the theorem, though the ``only if'' asserts the
existence of some consistent family with no extension and not the
inexactness of one particular fixed point. And the closure under intersection that the region
graph performs, rule~2 of Sec.~\ref{sec:regions}, is the construction that
turns a family of complexes into one on which consistency can be stated at
all: without the intersections as regions, ``agree on the overlap'' is not a
constraint the messages carry.

\paragraph{What is and is not in the theorem.} Two cautions, so that the
attribution is not read as more than it is. The theorem is about
\emph{existence} of an extension, not about the recursion finding it; on an
acyclic family the fixed point is unique and the recursion reaches it, but
that is \citeauthor{lauritzen1988}'s half, not \citeauthor{vorobev1962}'s.
And the theorem is an instance statement, like the GYO residue of
Sec.~\ref{sec:residue}: it says nothing about ensembles, and the local
criterion of Chapter~\ref{ch:exactness} --- the fraction $\phi$ of vertices
on short chordless cycles, Eq.~\eqref{eq:phishort} --- has no counterpart in
it. What the chygraph adds to the marginal problem is exactly that
ensemble: Sec.~\ref{sec:mg-nested} returns to it.

\section{Sheaves and the obstruction}
\label{sec:mg-sheaf}
\index{sheaf|textbf}
\index{contextuality|textbf}
\index{cohomology}

When the answer to the marginal problem is no, something is left over, and
\citet{abramsky2011} said what kind of thing it is. Their setting is
quantum contextuality --- measurement scenarios in which the outcome
statistics of each set of jointly measurable observables are given, and the
question is whether one hidden joint law reproduces them all --- and the
question is the marginal problem word for word, with the hypergraph of
jointly measurable sets in place of $H$. Their reading: the local laws form
a \emph{presheaf} over the hypergraph, assigning to each hyperedge its law
and to each inclusion $e'\subset e$ the marginalisation map; a global law is
a \emph{global section} of that presheaf; and the non-existence of one is a
failure of the presheaf to be a sheaf, which is the kind of failure
cohomology measures. \citet{abramsky2012} make the measurement: an
obstruction class in the first \v{C}ech cohomology of the presheaf, which
when non-zero certifies that the family is contextual in the strong sense.
The implication runs one way only: there are strongly contextual families
whose class vanishes, their Kochen--Specker example among them, so a
vanishing class settles nothing. \citet{abramsky2015} compute the class on
the standard paradoxes.

\paragraph{Part IV's error as a class.} This is the language for what
happens under overlap, with a caveat that governs the paragraph: whether
the fixed points of Part~IV glue has not been checked, and the two
diagnostics need not coincide, since a family can glue and still be wrong.
Chapter~\ref{ch:overlap} measured the error of the chygraph recursion on
two triangles sharing an edge, Table~\ref{tab:twotriangles}, and on a ring
of triangles glued at single atoms, Fig.~\ref{fig:ring} and
Table~\ref{tab:ring}, as a number: the discrepancy in $\ln Z$ or in a
magnetisation. In the presheaf reading the question asked of a fixed point
is a different one from the question the number answers: not how far the
beliefs are from the true marginals, but whether they are the marginals of
\emph{any} law at all. The object that answers, when it is non-zero, is a
cohomology class supported on the cycle. On the two triangles the cycle is
the four-cycle through their outer vertices, whose chord the recursion
cannot see because the chord's edge belongs to both complexes and is
counted twice; on the ring it is the ring itself. Two things should be
said about the class that the number does not carry. It vanishes for every
consistent family when the hypergraph is $\alpha$-acyclic, which covers
every instance Chapter~\ref{ch:gbp} found exact; on a cyclic hypergraph it
can be non-zero, and Part~IV's instances are where to look for it, with the
one-way implication above in mind. And its support lies on the cycles over
complexes, so the number of independent such cycles, which
Chapter~\ref{ch:mobius} counts homologically and expands the loop series
in, bounds where the obstruction can live; it is not the dimension of the
cohomology group, which is that of a presheaf of formal sums of local
sections and not of the nerve.

The repairs of Part~IV read the same way. Merging the complexes that share
two atoms, Sec.~\ref{sec:merging}, replaces the two hyperedges on which the
class lives by one, on which there is no cycle to support it; the region
graph of Sec.~\ref{sec:regions} adds the intersection as a region, which is
a refinement of the cover, and a class that does not survive refinement was
never a global obstruction. What Chapter~\ref{ch:mobius} takes up is the
question of which classes survive: the M\"obius counting numbers of
Eq.~\eqref{eq:mobius} are the Euler characteristic of the cover, and the
first Betti number of the incidence structure is the number of independent
classes there can be.

\section{Marginals of marginals}
\label{sec:mg-nested}
\index{marginal problem!nested}
\index{hierarchical hypergraph}

\citeauthor{vorobev1962} and \citeauthor{kellerer1964} pose the marginal
problem on a hypergraph of one level: variables, and sets of variables. The
chygraph poses it on a hierarchy. A complex that contains complexes has, at
its fixed point, a law on its own state and on the states of its members,
which are themselves complexes with laws on \emph{their} members; the
consistency condition is now that a complex's law on a member agrees with
that member's own law on itself, at every level, and the question is whether
such a nested family glues to one law on all the atoms and all the complex
states.

\paragraph{The level-wise statement.} The natural statement is

\begin{center}
\emph{a consistent nested family of marginals glues to a global law\\
if and only if the flattened family $\{\{a\}\cup\del a\}$ is
$\alpha$-acyclic,}
\end{center}

the flattened family being the hyperedges of the factor graph of rule~3,
one per complex, with the complex's own state kept as a variable. It is
not a conjecture. The join-tree argument uses nothing about levels, so
Vorob'ev and Kellerer apply to the flattened family as it stands, and
what remained was a definition check: whether ``consistent'' as defined
level by level --- a container's law on a member complex agrees with that
complex's own law on itself, and laws agree on shared atoms --- is the
same as consistency on the flattened family. It is worth being precise
about which acyclicity, because the weaker one fails: with ``the factor
graph of $\alpha+I$ is a tree'' in place of $\alpha$-acyclicity the
``only if'' is false, two triangles sharing an edge with their states kept
having an $\alpha$-acyclic flattened family, so that every consistent
family glues, and a factor graph with a cycle.

\paragraph{The check, done.}\index{consistency!level-wise|textbf}\index{Vorob'ev--Kellerer theorem|textbf} Two hyperedges of the flattened family,
$\{a\}\cup\del a$ and $\{b\}\cup\del b$, overlap in one of two ways up to
symmetry. If $b\in\del a$ the overlap is $\{b\}\cup(\del a\cap\del b)$: the member's
own state, together with anything $a$ holds that is also a member of $b$.
If neither contains the other the overlap is $\del a\cap\del b$, the
members they share. Level-wise consistency is pairwise agreement on
single vertices --- the state $\{b\}$, or one shared member at a time ---
so the two notions coincide exactly when every pairwise overlap of the
flattened family is a single vertex, and differ as soon as one has two or
more: two containers sharing two members, or a container holding a
member together with one of that member's members. Where they differ the
level-wise notion is the weaker one, and a level-wise consistent family
need not be consistent, for the same reason as on two triangles sharing
an edge: agreement on $\{2\}$ and on $\{3\}$ says nothing about the law on
$(\sigma_{2},\sigma_{3})$, and the two complexes can carry opposite
correlations on the pair and still agree atom by atom, which is the
disagreement Sec.~\ref{sec:mg-marginal} measured. So the statement is:

\begin{center}
\emph{a nested family of marginals that agrees on every pairwise overlap
of the flattened family glues to a global law if and only if that family
is $\alpha$-acyclic; and level-wise agreement is the same hypothesis
exactly when every such overlap is a single vertex.}
\end{center}

The first clause is Vorob'ev and Kellerer verbatim; the second is the
check. For a strictly layered chygraph, where a container's members all
sit one level down, the single-vertex condition says that no two
containers share two members and no two complexes share two atoms, which
is Part~IV's pairwise condition carried up the levels; and where it fails
the remedy is the one Chapter~\ref{ch:gbp} applied at the bottom level,
to take the overlaps as regions and demand agreement on them jointly.
Section~\ref{sec:nested} is the one place in the book where a flattening
is done, and there it costs: CNF flattening of a nested clause is
\emph{not} the flattening of rule~3, because it eliminates the clause's
state rather than keeping it as a variable, and
Sec.~\ref{sec:flatteningcost} measured what that costs. The statement
above says that the rule-3 flattening, which keeps the state, costs
nothing.

\paragraph{Contextuality on a hierarchy.} The same move in
\citeauthor{abramsky2011}'s setting gives nested measurement scenarios:
contexts that contain contexts, a joint measurement whose outcome is itself
one of the observables of a larger joint measurement. Nothing in the
sheaf-theoretic account forbids that, and the obstruction class would live
on the factor graph of $\alpha+I$ as before. What the chygraph then offers
is an \emph{ensemble} of scenarios: a layered random chygraph is a random
measurement scenario, the question ``does a random scenario admit a global
section'' is the question ``is the flattened family $\alpha$-acyclic'', and
Chapter~\ref{ch:exactness}'s acyclicity percolation is a contextuality
threshold --- below it a random scenario is non-contextual with probability
tending to one, above it a finite fraction of its contexts sit on an
obstruction. That is a reading and not a result; no scenario has been built
to test it.

\paragraph{A possible prior name.} One more, to be verified before it is
relied on, because it may already exist under another name. The
\emph{hierarchical hypergraphs} that \citet{alvarezpicallo2022} introduce
for rewriting string diagrams in monoidal closed categories are hypergraphs
whose hyperedges may contain hypergraphs, to any depth, with the containing
edge acting as an interface to the contained one. That is close to a
chygraph: a complex is a hyperedge, its members are the vertices and
hyperedges it contains, and the interface is the complex's own state. One
difference is visible from the definitions as far as they can be read from
here: their nesting appears to be strict, a hyperedge owning the hypergraph
it contains, whereas a chygraph lets an atom lie in several complexes of
the same level, which is the whole of Part~IV. Whether the two coincide on
the strictly nested case, or which is the more general, has not been
checked. If they do, each side has something for the other.
The chygraph supplies that community with random instances --- the
ensembles of Sec.~\ref{sec:ensembles} --- and with a statistical mechanics
of string diagrams, thresholds and free energies read off the four moment
matrices of Eq.~\eqref{eq:moments}. And they supply the book with a
rewriting theory for the merge closure of Chapter~\ref{ch:metacomplex}: the
closure of Sec.~\ref{sec:merging} is a rewriting system on the chygraph,
two complexes sharing two atoms rewritten to their union and the process run
to a fixed point, and whether that fixed point is unique and reached in a
bounded number of rounds is exactly the confluence and termination that
rewriting theory decides. Section~\ref{sec:merging} counts the rounds; a rewriting argument would
bound them.

\section{What it gives the book}
\label{sec:mg-gives}

Three things.

\emph{The exactness theorem with its proper attribution.} The chain of
Sec.~\ref{sec:threenames} is the hypothesis of a theorem of
\citet{vorobev1962} and \citet{kellerer1964}, and the exactness of the
recursion on it is that theorem's conclusion, with \citet{lauritzen1988}
supplying the algorithm. The book's contribution to the instance statement
is the certificate --- leaf removal on the incidence structure, which is
Chapter~\ref{ch:cover}'s tool --- and its contribution to the theory is the
ensemble version, which the theorem does not have.

\emph{The treelike assumption as a definition.} In a Markov category
``the complexes at an atom carry independent information given the atom''
is a factorisation of a morphism, the same for every model in
Table~\ref{tab:ce-substitutions}, so what the book has been calling an
assumption about the structure is exactly that and not an assumption about
the Hamiltonian. Section~\ref{sec:assumes} said so in words; the category
says it in a way that cannot be misread.

\emph{Part IV's error as a class.} The number Chapter~\ref{ch:overlap}
measures is the size of an obstruction whose support is the cycle space
over complexes and whose vanishing is $\alpha$-acyclicity. That is the
reading Chapter~\ref{ch:mobius} needs in order to organise the corrections
by loop rather than by instance.

\section{What the field could get}
\label{sec:mg-field}

\emph{A hierarchy.} The marginal problem has been posed on one level for
sixty years. The chygraph poses it on a hierarchy, and
Sec.~\ref{sec:mg-nested} states the level-wise theorem as a corollary of
the one-level one, the definition check done: level-wise agreement is the
theorem's hypothesis exactly when the flattened family's overlaps are
single vertices.

\emph{An ensemble.} Contextuality is studied one scenario at a time. A
layered random chygraph is a random scenario, and the acyclicity transition
of Chapter~\ref{ch:exactness} says at what density a random scenario
acquires an obstruction. Whether the physically interesting scenarios sit
on one side of that transition is a question the ensemble can be asked and
a single scenario cannot.

\emph{Random string diagrams.} If hierarchical hypergraphs are chygraphs,
then everything in Parts~II and~III is a statistical mechanics of string
diagrams --- percolation of a random diagram, its free energy, its
thresholds --- and Part~IV is the theory of when a diagram can be rewritten
to a tree. The merge closure is a rewriting system whose confluence and
termination have been measured and not proved; a rewriting theory would
prove them, or say on which chygraphs they fail.

\section{Checks}
\label{sec:mg-checks}
\index{consistency!of clique beliefs, measured}

\checkhead{Known results}
The Vorob'ev--Kellerer theorem is quoted, not re-proved; its two directions
are illustrated on a tree and a triangle of pairwise marginals in
Sec.~\ref{sec:mg-marginal}, and its hypothesis is identified with the
chain of Sec.~\ref{sec:threenames} by citation to \citet{beeri1983}.

\checkhead{New results}
The two measurements of Sec.~\ref{sec:mg-marginal}.
\code{statmech/\allowbreak probe/\allowbreak pair\_\allowbreak consistency.\allowbreak py}
takes the $240$ runs of Chapter~\ref{ch:overlap} (Karrer--Newman and real
instances as there, hyperbolic redrawn, Sec.~\ref{sec:cavityfail}) at the
stable fixed point of the chygraph recursion, with the double count and
with each bond in one clique, and for every pair of cliques sharing two or
more atoms compares the two clique beliefs marginalised to the shared
atoms, and to each atom singly. Single atoms agree to $3\times10^{-13}$
throughout; pairs differ by up to $0.149$ and $0.328$. The region-graph
side reads the consistency the runs of Chapter~\ref{ch:gbp} recorded, the
largest violation of parent-to-child marginalisation over the region
graph's edges, which on a family closed under intersection is agreement on
every overlap: on the $186$ of $240$ runs whose iteration settled it is
below $10^{-12}$ on the chordal instances and below $8.5\times10^{-8}$ on
all but one non-chordal instance (that one at $0.057$). The $54$ runs that
did not settle are not consistent and are not quoted. The definition
check of Sec.~\ref{sec:mg-nested} is a case analysis and needs no
computation.

\checkhead{Partial results}
Where the region-graph iteration did not settle nothing is claimed. The
magnetisations recorded for the Karrer--Newman and real runs of
Chapter~\ref{ch:gbp} are zero on most of them, so those fixed points are
the symmetric one, which on a junction tree is the exact answer and off
it is a fixed point whose stability was not checked; Sec.~\ref{sec:cavityfail}
found the same point unstable for the chygraph recursion on most of the
same instances.
